% !PRINTED none
% Part-title (Abschnitt-opening) page: "Dreizehnter Abschnitt. Äquivalenz nach
% Zahlgruppen." centred, rule, then the opening of § 98 "Zahlgruppen in den
% Ordnungen." with running text below. No page number is printed in the
% running head area on this page.

\begin{center}
\gesperrt{Dreizehnter Abschnitt.}

Äquivalenz nach Zahlgruppen.
\end{center}

\rule{0.3\textwidth}{0.4pt}

\begin{center}
\textbf{§ 98. \gesperrt{Zahlgruppen in den Ordnungen.}}
\end{center}

Um die Beziehungen der quadratischen Irrationalzahlen zu
den Ordnungen der quadratischen Formen genauer zu erforschen,
leiten wir aus den Zahlen des Körpers $\Omega$ \gesperrt{Zahlgruppen} nach
folgenden Gesichtspunkten ab.

1. Wir nehmen eine beliebige natürliche Zahl $Q$ und schließen
von den ganzen Zahlen in $\Omega$ zunächst alle die aus, die nicht
teilerfremd zu $Q$ sind. Von den übrigen nehmen wir auch noch
beliebige Quotienten.

Die so erhaltenen ganzen und gebrochenen Zahlen $\eta$ nennen
wir \gesperrt{fremd gegen $Q$}. Ihre Gesamtheit sei mit $O$ bezeichnet. Sie
bilden insofern eine \gesperrt{Gruppe}, als Multiplikation und Division
zweier dieser Zahlen immer Zahlen desselben Systems liefern.

2. Wir verfahren nun ebenso mit den Zahlen der Ordnung $[Q]$,
d.\,h. wir schließen alle Zahlen dieser Ordnung, die mit $Q$ nicht
relativ prim sind, aus und bilden auch noch die Quotienten
beliebiger zweier der übrig gebliebenen Zahlen. Wir bezeichnen
diese Zahlen in $\eta'$ % [sic] Gedruckt ist in; grammatisch wäre mit zu erwarten. und ihre Gesamtheit mit $O'$, und nennen sie
ganze und gebrochene Zahlen der Ordnung $[Q]$. Die Zahlen $O'$
sind alle in $O$ enthalten und bilden gleichfalls der Multiplikation
und Division gegenüber eine \gesperrt{Gruppe}. Wir nennen sie die
\gesperrt{Gruppe der Ordnung $[Q]$} oder kurz eine \gesperrt{Ordnungsgruppe}.

Jede Zahl $\eta$ in $O$ kann durch Multiplikation mit einer ganzen
rationalen zu $Q$ \gesperrt{teilerfremden} Zahl $n$ in eine ganze Zahl $\omega$
verwandelt werden.

Wenn $\eta'$ eine Zahl in $O'$ ist, so kann man den Nenner
rational machen und findet, daß $n\eta'$ eine ganze Zahl der Ordnung
$[Q]$ und folglich mit einer rationalen Zahl nach dem Modul $Q$
kongruent ist.
